% Corrigé intégré automatiquement pour la banque Exercices.
% Source : B2_ProposerModele/B2_16_Hyperstatisme/84_Nacelle/84_Nacelle.tex
\begin{corrigebox}[Corrigé]
\CorrigeQuestion{1}
On recense les inconnues de liaison réellement indépendantes et le nombre d'équations d'équilibre disponibles après prise en compte des mobilités. Le degré d'hyperstatisme est obtenu par le bilan $h=N_s-r_s$, puis contrôlé en identifiant les inconnues redondantes.
\CorrigeQuestion{2}
On recense les inconnues de liaison réellement indépendantes et le nombre d'équations d'équilibre disponibles après prise en compte des mobilités. Le degré d'hyperstatisme est obtenu par le bilan $h=N_s-r_s$, puis contrôlé en identifiant les inconnues redondantes.
\CorrigeQuestion{3}
On identifie les classes d'équivalence cinématique, les surfaces de contact et les mouvements relatifs autorisés. Chaque contact est remplacé par la liaison normalisée correspondante, puis les axes et points caractéristiques sont reportés sur le schéma cinématique minimal.
\CorrigeQuestion{4}
La résolution consiste à identifier les données et l'inconnue, choisir la loi du cours adaptée, écrire l'équation symbolique avant toute application numérique, puis vérifier l'homogénéité, le signe et l'ordre de grandeur du résultat. Les hypothèses utilisées doivent être explicitement contrôlées à la fin.
\CorrigeQuestion{5}
La résolution consiste à identifier les données et l'inconnue, choisir la loi du cours adaptée, écrire l'équation symbolique avant toute application numérique, puis vérifier l'homogénéité, le signe et l'ordre de grandeur du résultat. Les hypothèses utilisées doivent être explicitement contrôlées à la fin.
\end{corrigebox}
